\documentclass[10pt]{amsart}

\usepackage[english]{babel}
\usepackage{amssymb}
\usepackage[leqno]{amsmath}
\usepackage{mathrsfs}
\usepackage{stmaryrd}
\usepackage{chemarrow}
\usepackage[norelsize]{algorithm2e}
\usepackage{enumerate}
\usepackage{graphicx}
\usepackage{subcaption}
\usepackage[all]{xy}
\usepackage{tikz}
\usetikzlibrary{arrows}
\usepackage{multirow}
\usepackage[colorlinks=true,allcolors=blue]{hyperref}
\usepackage{booktabs}
\usepackage{float}
\usepackage{placeins}

\newtheorem{theorem}{Theorem}[section]
\newtheorem{lemma}[theorem]{Lemma}
\newtheorem{corollary}[theorem]{Corollary}
\newtheorem{proposition}[theorem]{Proposition}
\newtheorem{definition}[theorem]{Definition}
\newtheorem{example}[theorem]{Example}
\newtheorem{exercise}[theorem]{Exercise}
\newtheorem{question}[theorem]{Question}
\newtheorem{remark}[theorem]{Remark}
\newtheorem{assumption}[theorem]{Assumption}
\newtheorem{alg}[theorem]{Algorithm}

\newcommand{\dx}{\,{\rm d}x}
\newcommand{\dd}{\,{\rm d}}
\newcommand{\bs}{\boldsymbol}
\newcommand{\mcal}{\mathcal}

\DeclareMathOperator*{\img}{img}
\newcommand{\sign}{\operatorname{sign}}
\newcommand{\curl}{\operatorname{curl}}
\renewcommand{\div}{\operatorname{div}}
\newcommand{\grad}{\operatorname{grad}}
\newcommand{\tr}{\operatorname{tr}}
\DeclareMathOperator*{\rot}{rot}
\DeclareMathOperator*{\var}{Var}
\newcommand{\dev}{\operatorname{dev}}
\newcommand{\sym}{\operatorname{sym}}
\newcommand{\skw}{\operatorname{skw}}
\newcommand{\spn}{\operatorname{spn}}
\newcommand{\mspn}{\operatorname{mspn}}
\newcommand{\mskw}{\operatorname{mskw}}
\newcommand{\vskw}{\operatorname{vskw}}
\newcommand{\vspn}{\operatorname{vspn}}
\newcommand{\hess}{\operatorname{hess}}

\numberwithin{equation}{section}

\begin{document}

\title[regularity of strain gradient elasticity]
{A uniform regularity of the strain gradient elasticity model on polytopes}

% \author{Xuehai Huang}
% \address{School of Mathematics, Shanghai University of Finance and Economics,
% Shanghai 200433, China}
% \email{huang.xuehai@sufe.edu.cn}

% \author{Xinyue Zhao}
% \address{School of Mathematics, Shanghai University of Finance and Economics,
% Shanghai 200433, China}
% \email{zhaoxinyue20210921@163.com}

% \thanks{The first author was supported by the National Natural Science
% Foundation of China Project 12671432.}


\begin{abstract}
A uniform regularity of the strain gradient elasticity model on polytopes is derived under reasonable assumptions.
\end{abstract}

% \keywords{
% weakly symmetric and traceless tensors, tangential--normal finite
% elements, Brinkman equations, distributional mixed finite element methods
% }

% \subjclass[2020]{65N30, 65N12, 65N22, 76D07, 76S05}

\maketitle

\section{Introduction}
Let $\Omega\subset \mathbb{R}^d$ ($d\geq2$) be a polytope. Consider the following strain gradient elasticity (SGE) model: 
\begin{equation}
\begin{cases}\label{eq:straingradproblem}
-\div\big((\boldsymbol I-\iota^2\Delta)\boldsymbol\sigma(\boldsymbol u)\big)=\boldsymbol f & \textrm{ in }  \Omega,\\
\boldsymbol u=\partial_{\boldsymbol n}\boldsymbol u=\boldsymbol0 & \textrm{ on } \partial\Omega,
\end{cases}
\end{equation}
where $\boldsymbol n$ is the unit outward normal to $\partial\Omega$, $\boldsymbol u=(u_1, \ldots, u_d)^{\intercal}$ is the displacement field, $\boldsymbol f$ is the applied force, $\boldsymbol\sigma(\boldsymbol u)=2\mu\boldsymbol\varepsilon(\boldsymbol u)+\lambda(\div\boldsymbol u)\boldsymbol I$, and $\boldsymbol\varepsilon(\boldsymbol u)=(\varepsilon_{ij}(\boldsymbol u))_{d\times d}$ is the strain tensor field with
$
\varepsilon_{ij}(\boldsymbol u)=1/2(\partial_i u_j+\partial_ju_i).
$
Here $\lambda$ and $\mu$ are the Lam\'e constants, and $\boldsymbol I$ is the identity tensor field. The resulting stress field is given by $\tilde{\boldsymbol\sigma}(\boldsymbol u)=(\boldsymbol I-\iota^2\Delta)\boldsymbol\sigma(\boldsymbol u)$, where $\iota\in (0,1)$ denotes the size parameter of the material under discussion.

In this paper, we establish the following regularity estimate for problem~\eqref{eq:straingradproblem}:
\begin{equation}\label{intro:straingradregularity}	
\begin{aligned}	
&|\boldsymbol u-\boldsymbol u_0|_{1} +\iota\|\boldsymbol u\|_2  + \iota^{2}\|\boldsymbol u\|_{3}\\
&\qquad+\lambda\|\div(\boldsymbol u-\boldsymbol u_0)\|_0 + \lambda\iota|\div\boldsymbol u|_1 + \lambda\iota^2\|\div\boldsymbol u\|_2 
\ \lesssim\ \iota^{1/2}\|\boldsymbol f\|_0,
\end{aligned}
\end{equation}
under Assumptions~\ref{RegularityAssumption0}--\ref{RegularityAssumption2}, via a standard perturbation argument.
The estimate~\eqref{intro:straingradregularity} was proved for smooth domains in~\cite{LiaoMingXu2023}, and we built up the regularity \eqref{intro:straingradregularity} for polytopes in \cite{ChenHuangHuang2023}.

A key ingredient in the proof of Lemma~3.7 in~\cite{ChenHuangHuang2023} is the following div-surjectivity statement:
\[
\div\big( H^3(\Omega;\mathbb R^d)\cap H_0^2(\Omega;\mathbb R^d)\big)
= H^2(\Omega)\cap H_0^1(\Omega)\cap L_0^2(\Omega),
\]
which is used there without proof and is not obvious. Indeed, the analogous surjectivity at one lower order fails:
\[
\div\big( H^2(\Omega;\mathbb R^2)\cap H_0^1(\Omega;\mathbb R^2)\big)
\neq H^1(\Omega)\cap L_0^2(\Omega).
\]
As pointed out in~\cite[Theorem~3.1]{ArnoldScottVogelius1988}, the space
$\div\big( H^2(\Omega;\mathbb R^2)\cap H_0^1(\Omega;\mathbb R^2)\big)$
is a proper closed subspace of $H^1(\Omega)\cap L_0^2(\Omega)$. In this paper, we prove in Lemma~\ref{lem:main} that the div-surjectivity
\[
\div\big( H^3(\Omega;\mathbb R^2)\cap H_0^2(\Omega;\mathbb R^2)\big)
= H^2(\Omega)\cap H_0^1(\Omega)\cap L_0^2(\Omega)
\]
indeed holds at least in two dimensions.


For later uses, we introduce some notation.
Given a bounded domain $D\subset \mathbb{R}^{d}$
and integers $m,k\geq0$, denote by $H^m(D)$ the standard Sobolev space on $D$
with norm $\|\cdot\|_{m,D}$ and semi-norm $|\cdot|_{m,D}$, and $H_0^m(D)$ the closure of $C_0^{\infty}(D)$ with respect to $\|\cdot\|_{m,D}$. We abbreviate $\|\cdot\|_{m,\Omega}$ and $|\cdot|_{m,\Omega}$ as $\|\cdot\|_{m}$ and $|\cdot|_{m}$ respectively for $D=\Omega$.
The notation $(\cdot,\,\cdot)_{D}$ symbolizes the standard $L^2$ inner product on $D$, and the subscript will be omitted when $D=\Omega$.
% Let $\mathbb{P}_{k}(D)$ be the set of all polynomials on $D$ with the total degree up to $k$. 
Set  $H^m(D; \mathbb R^d) :=H^m(D)\otimes\mathbb R^d$, $H_0^m(D; \mathbb R^d) :=H_0^m(D)\otimes\mathbb R^d$.
%  and $\mathbb{P}_k(D; \mathbb R^s) :=\mathbb{P}_k(D)\otimes\mathbb R^s$. 
Let $L^2_0(D)$ be the space of functions in $L^2(D)$ with vanishing integral average values.
%   \begin{align*}
% & H^m(D; \mathbb R^s) = \{\boldsymbol v=(v_1,v_2,\cdots,v_s)^T| v_i\in H^m(D),1\leq i \leq s \}, \\
% & H^m_0(D; \mathbb R^s) = \{\boldsymbol v\in H^m(D; \mathbb R^s)| \boldsymbol v = \frac {\partial\boldsymbol v}{\partial \boldsymbol n} = \cdots = \frac {\partial^{m-1}\boldsymbol v}{\partial \boldsymbol n^{m-1}}=0 \text{ on } \partial D \}, \\
% & \mathbb{P}_k(D; \mathbb R^s) = \{\boldsymbol v= [v_1,v_2,\cdots,v_s]^T| v_i\in \mathbb{P}_k(D),1\leq i \leq s \}, \\
% &L^2_0(D) = \{p\in L^2(D) | \int_{D} p {\rm d}x=0 \}.
% \end{align*}

For a scalar function $w$ and a vector valued function $\boldsymbol v = (v_1,v_2)^{\intercal}$ in two dimensions, introduce two usual differential operations
$\curl w = (\partial_2 w, - \partial_1 w)^{\intercal}$. 
% and $\mathbb Rot \boldsymbol v = \partial_1 v_2 - \partial_2 v_1.$

% In the rest of this subsection, we present some basic inequalities for later theoretical analysis.
As usual, we use $\lesssim$ to represent $\leq C$, where $C$ may be a generic positive constant independent of %the mesh size $h$, 
the material parameter $\iota$ and the Lam\'e coefficient $\lambda$. And $a\eqsim b$ indicates $a\lesssim b\lesssim a$. 
% Assume $\Omega$ is a star-shaped domain \cite{Brenner2008FEMtheory}.
% We have the multiplicative trace inequality in the form \cite[Theorem 1.5.1.10]{Grisvard1985}
% \begin{equation}\label{eq:multiplicativetraceineqty}
% \|v\|_{\partial D}^2 \lesssim \|v\|_{0,D} \|v\|_{1,D} \quad\forall~v\in H^1(D),
% \end{equation}
% which implies readily the scaled trace inequality %\cite[(2.18)]{BrennerSung2018}:
% \begin{equation}\label{eq:traceineqty}
% \|v\|_{0,\partial D}\lesssim h_D^{-1/2}\|v\|_{0,D} + h_D^{1/2}|v|_{1,D} \quad\forall~v\in H^1(D).
% \end{equation}

% \begin{equation}\label{eq:korninequalityH2}
% ( 1-1/\sqrt{2})\left(|\boldsymbol v|_1^2 + |\boldsymbol v|_2^2\mathbb Right)\lesssim \|\boldsymbol\varepsilon(\boldsymbol v)\|_{0}^2 + \|\nabla\boldsymbol\varepsilon(\boldsymbol v)\|_0^2 \quad\forall~\boldsymbol v\in H^1_0(\Omega;\mathbb{R}^d)\cap H^2(\Omega;\mathbb{R}^d).
% \end{equation}

% The following inequality can be found in \cite[Theorem 2.2]{GiraultRaviart1986}:
% \begin{equation}\label{eq:poincareL2m1}
% \|p\|_0\lesssim \|p\|_{-1}+\|\nabla p\|_{-1}\quad\forall~p\in L^2(\Omega).
% \end{equation}

% Recall the continuity of the right inverse of the divergence operator in \cite{CostabelMcIntosh2010}.
% \begin{lemma}\label{lm:divrightinverse}
% For $p\in H_0^{m}(\Omega)\cap L_0^2(\Omega)$ with non-negative integer $m$, there exists $\boldsymbol{v}\in H_0^{m+1}(\Omega;\mathbb R^d)$ such that $\div\boldsymbol{v}=p$ and
% \begin{equation*}
% \|\nabla\boldsymbol v\|_{j}\lesssim \|p\|_{j} \quad {\rm  for }\;  j=0,1,\ldots, m.
% \end{equation*}
% \end{lemma}
% \cref{lm:divrightinverse}


% \subsection{Mixed formulation}
% %To construct a discrete method for \eqref{eq:straingradproblem} possessing uniform convergence with respected to $\lambda$ and $\iota$, Referring to \cite{LiaoMingXu2021taylorhood},We mimic Lam\'e system by introducing $p=\lambda\div\boldsymbol u$
% In order to solve the problem \eqref{eq:straingradproblem} with a mixed finite element method, we recall its mixed formulation in \cite{LiaoMingXu2021taylorhood}.
% % By introducing an auxiliary variable $p=\lambda\div\boldsymbol u$ as for the Lam\'e system, one can reformulate \eqref{eq:straingradproblem} in the form \cite{LiaoMingXu2021taylorhood}
% % \begin{equation*}
% % \begin{cases}
% % -\div\big((I-\iota^2\Delta)(2\mu\boldsymbol\varepsilon(\boldsymbol u))\big) -\nabla\big((I-\iota^2\Delta)p\big)=\boldsymbol f & \textrm{ in } \Omega,\\
% % \div\boldsymbol u -\dfrac{1}{\lambda}p=0 & \textrm{ in } \Omega,\\
% % \boldsymbol u=\partial_{\boldsymbol{n}}\boldsymbol u=\boldsymbol 0 & \textrm{ on } \partial\Omega.
% % \end{cases}
% % \end{equation*}
% % Then, by integration by parts, we can deduce that
% The variational formulation of \eqref{straingradLameproblem} is to find $\boldsymbol u\in V:=H_0^2(\Omega; \mathbb R^d)$ and $p\in Q:=H_0^1(\Omega)\cap L^2_0(\Omega)$ such that
% \begin{equation}\label{eq:straingradmixedform}
% \begin{cases}
% a(\boldsymbol u, \boldsymbol v)
% +b(\boldsymbol v, p)=(\boldsymbol f, \boldsymbol v)&  \forall~\boldsymbol v\in V,\\
% b(\boldsymbol u, q)-c(p,q)=0 & \forall~q\in Q,
% \end{cases}
% \end{equation}
% where
% % \[
% % a(\boldsymbol u, \boldsymbol v)=2\mu(\boldsymbol\varepsilon(\boldsymbol u), \boldsymbol\varepsilon(\boldsymbol v))_{\iota}, b(\boldsymbol v, p)=(\div\boldsymbol v, p)_{\iota}, c(p,q)=(p, q)_{\iota}/\lambda,
% % \]
% \begin{align*}
% a(\boldsymbol u, \boldsymbol v)&:=2\mu(\boldsymbol\varepsilon(\boldsymbol u), \boldsymbol\varepsilon(\boldsymbol v))_{\iota}=2\mu\big((\boldsymbol\varepsilon(\boldsymbol u), \boldsymbol\varepsilon(\boldsymbol v))+\iota^2(\nabla\boldsymbol\varepsilon(\boldsymbol u), \nabla\boldsymbol\varepsilon(\boldsymbol v))\big),	 \\
% b(\boldsymbol v, p)&:=(\div\boldsymbol v, p)_{\iota}=(\div\boldsymbol v, p)+\iota^2(\nabla\div\boldsymbol v, \nabla p), \\
% c(p,q)&:=(p, q)_{\iota}/\lambda=\big((p, q)+\iota^2(\nabla p, \nabla q)\big)/\lambda
% \end{align*}
% with a weighted $H^1$ inner product $(p, q)_{\iota}:=(p, q)+\iota^2(\nabla p, \nabla q)$.
% To make the forthcoming discussion in a compact way, we further introduce the following weighted Sobolev norms over $V$ and $Q$:
% \[
% \|\boldsymbol v\|_V : = \left(|\boldsymbol v|_1^2 + \iota ^2 |\boldsymbol v|_2^2\mathbb Right)^{1/2}\quad \forall\; \boldsymbol v\in V; \quad \|q\|_Q : = \left(\|q\|_0^2 + \iota ^2 |q|_1^2\mathbb Right)^{1/2}\quad \forall\; q\in Q.
% \]

% Now, it is easy to check that the linear forms $a(\cdot,\cdot)$, $b(\cdot,\cdot)$ and $c(\cdot,\cdot)$ are bounded on $V\times V$, $V\times Q$ and $Q\times Q$, respectively. The coercivity of the linear form $a(\cdot,\cdot)$ on $V\times V$ results from Korn's inequalities~\eqref{eq:korninequality1st}-\eqref{eq:korninequalityH2}. And the inf-sup condition
% \begin{equation*}	
% \|q\|_{Q}\lesssim  \sup\limits_{\boldsymbol v \in V}\frac {b(\boldsymbol v,q)} {\|\boldsymbol v\|_V } \quad \forall~q\in Q
% \end{equation*}
% holds from Lemma~\ref{lm:divrightinverse}. Applying the Babu{\v{s}}ka-Brezzi theory \cite{BoffiBrezziFortin2013}, we have the stability
% % $$
% % \|\boldsymbol u\|_V+\|p\|_Q \lesssim \sup_{\boldsymbol v\in V, q\in Q}\frac{a(\boldsymbol u, \boldsymbol v)+b(\boldsymbol v, p)-b(\boldsymbol u, q)}{\|\boldsymbol v\|_V+\|q\|_Q} \quad\forall~\boldsymbol u\in V, p\in Q,
% % $$
% % which immediately gives
% \begin{equation*}	
% \|\boldsymbol u\|_V+\|p\|_Q \lesssim \sup_{\boldsymbol v\in V,\, q\in Q}\frac{a(\boldsymbol u, \boldsymbol v)+b(\boldsymbol v, p)-b(\boldsymbol u, q)+c(p,q)}{\|\boldsymbol v\|_V+\|q\|_Q} \quad\forall~\boldsymbol u\in V, p\in Q.
% \end{equation*}
% Therefore, the mixed formulation \eqref{eq:straingradmixedform} is well-posed.
% % with the stability bound
% % $$
% % \|\boldsymbol u\|_V+\|p\|_Q \lesssim \|\boldsymbol f\|_{V'}=\sup_{\boldsymbol v\in V}\frac{(\boldsymbol f, \boldsymbol v)}{\|\boldsymbol v\|_V}.
% % $$
% % \begin{align}
% % % \nonumber to remove numbering (before each equation)
% %   a(\boldsymbol v,\boldsymbol v) \gtrsim \|\boldsymbol v\|_{V} \quad \forall \boldsymbol v\in V, &\qquad c(q,q)\geq \|q\|_{Q} \quad \forall q\in Q,  \notag \\
% %   \|q\|_{Q}\lesssim  &\sup\limits_{\boldsymbol v \in V}\frac {b(\boldsymbol v,q)} {\|\boldsymbol v\|_V } \quad \forall q\in Q. \label{infsupcondition}
% % \end{align}

% %  By the Babu\v{s}ka-Brezzi theory, problem \eqref{eq:straingradmixedform} is well-posed if and only if
% %  \begin{eqnarray}
% %  % \nonumber to remove numbering (before each equation)
% %    \sqrt{ a( \boldsymbol v, \boldsymbol v ) } + \sup \limits_{p\in Q} \frac {b(\boldsymbol v, q )} {\|q\|_Q}   & \eqsim & \|\boldsymbol v\|_V,\quad \boldsymbol v\in V, \label{eq:wellposed1}\\
% %    \sqrt{ c( q, q ) } + \sup \limits_{v\in V} \frac {b(\boldsymbol v, q )} {\|\boldsymbol v\|_V}  & \eqsim &  \| q\|_P \label{eq:wellposed2}\quad q\in Q.
% %  \end{eqnarray}
% %   Clearly we have \eqref{eq:wellposed1} by Lemma \ref{lm:korn}
% %   and \eqref{eq:wellposed2} follows from the inf-sup condition \eqref{infsupcondition} immediately.


The rest of this paper is organized as follows. Three assumptions are presented in Sect.~\ref{sec:assumption}.
In Sect.~\ref{sec:regularity} we establish a uniform regularity for the SGE model under previous assumptions.


\section{Regularity assumptions}\label{sec:assumption}

% In order to derive the robust error estimates for the proposed mixed finite element method in the next section, we require to develop a series of regularity estimates for the underlying (system of) partial differential equations. In fact, these results are interesting themselves.
This section is devoted to presenting some regularity assumptions, and prove the regularity assumptions in two dimensions.

% \subsection{Assumptions}
\begin{assumption}\label{RegularityAssumption0}
Let $f \in L_0^2(\Omega)$, and $u\in H_0^1(\Omega)$ be the solution of the Poisson equation $-\Delta u=f$. Assume $u\in H^2(\Omega)$ admits the estimate
\begin{equation}\label{eq:poissonH2regularity}
\|u\|_2\lesssim \|f\|_{0}.
\end{equation}
\end{assumption}
The $H^2$ regularity \eqref{eq:poissonH2regularity} of the Poisson equation is well-known for convex domains \cite[Theorem 3.2.1.2]{Grisvard1985}.


Consider the following fourth order Stokes problem: Find $(\boldsymbol\phi, p)\in H_0^2(\Omega; \mathbb R^d)\times(H_0^1(\Omega)\cap L_0^2(\Omega))$ such that
\begin{equation}\label{eq:fourthorderStokes}
\begin{cases}
\Delta^2\boldsymbol\phi + \nabla\Delta p=\boldsymbol f & {\rm in }\; \Omega,\\
\Delta\div\boldsymbol\phi= g & {\rm in }\; \Omega,\\
\boldsymbol\phi=\partial_{\boldsymbol{n}}\boldsymbol\phi=\boldsymbol 0 & {\rm on }\; \partial\Omega,
\end{cases}
\end{equation}
where $\boldsymbol f \in H^{-2}(\Omega; \mathbb{R}^d)$ and $g\in H^{-1}(\Omega)$.
The well-posedness of the fourth order Stokes problem \eqref{eq:fourthorderStokes} follows from the div surjection  (cf. \cite[Theorem 1.1]{CostabelMcIntosh2010})
$$
\div H_0^2(\Omega; \mathbb R^d) = H_0^1(\Omega)\cap L_0^2(\Omega).
$$
The proof of the well-posedness of the fourth order Stokes problem \eqref{eq:fourthorderStokes} with $g=0$ in \cite[Lemma 3.1]{ChenHuangHuang2023} can be easily adapted to the general $g\in H^{-1}(\Omega)$. And it holds the stability estimate
\begin{equation}\label{eq:fourthorderStokesregularity0}
\|\boldsymbol\phi\|_2+\|p\|_1\lesssim \|\boldsymbol f\|_{-2} + \|g\|_{-1}.
\end{equation}


We will recall the two assumptions in \cite{ChenHuangHuang2023} with a slight modification.
Let $\chi\in H_0^1(\Omega)\cap H^2(\Omega)$ be the function solving
\[
-\Delta\chi=1\quad\text{in }\Omega,\qquad \chi=0\quad\text{on }\partial\Omega.
\]
Define the closed codimension--one subspace
\[
L^2_\chi(\Omega):=\big\{ g\in L^2(\Omega)\,:\, (g,\chi)_{L^2(\Omega)}=0\big\}.
\]

\begin{lemma}\label{lem:L2chi}
It holds that
$$
\Delta\div\big(H^3(\Omega;\mathbb R^d)\cap H_0^2(\Omega;\mathbb R^d)\big)\subseteq \Delta(H^2(\Omega)\cap H_0^1(\Omega)\cap L_0^2(\Omega)) \subseteq L^2_\chi(\Omega).
$$
\end{lemma}
\begin{proof}
By $\div\big(H^3(\Omega;\mathbb R^d)\cap H_0^2(\Omega;\mathbb R^d)\big)\subseteq (H^2(\Omega)\cap H_0^1(\Omega)\cap L_0^2(\Omega))$, it follows
$$
\Delta\div\big(H^3(\Omega;\mathbb R^d)\cap H_0^2(\Omega;\mathbb R^d)\big)\subseteq \Delta(H^2(\Omega)\cap H_0^1(\Omega)\cap L_0^2(\Omega)).
$$	
For any $p\in (H^2(\Omega)\cap H_0^1(\Omega)\cap L_0^2(\Omega))$, we have
\begin{align*}
(\Delta p,\chi)=(p,\Delta\chi)=-(p,1)=0,
\end{align*}
which implies $\Delta(H^2(\Omega)\cap H_0^1(\Omega)\cap L_0^2(\Omega)) \subseteq L^2_\chi(\Omega)$.
\end{proof}

\begin{lemma}
Assume the $H^2$ regularity \eqref{eq:poissonH2regularity} of the Poisson equation holds.
Then
\begin{equation}\label{eq:Deltasurjection}
\Delta(H^2(\Omega)\cap H_0^1(\Omega)\cap L_0^2(\Omega)) = L^2_\chi(\Omega).
\end{equation}
\end{lemma}
\begin{proof}
For any $g\in L^2_\chi(\Omega)$, let $v=\Delta^{-1}g\in H_0^1(\Omega)$.
By the $H^2$ regularity \eqref{eq:poissonH2regularity}, we have $v\in H^2(\Omega)\cap H_0^1(\Omega)$. Apply the integration by parts to get
$$
(v,1)=-(v,\Delta\chi)=-(\Delta v,\chi)=-(g,\chi)=0,
$$
which implies $v\in L_0^2(\Omega)$ and thus $g=\Delta v\in \Delta(H^2(\Omega)\cap H_0^1(\Omega)\cap L_0^2(\Omega))$.
\end{proof}

The two assumptions are presented as follows.
\begin{assumption}\label{RegularityAssumption1}
Let $\boldsymbol f \in H^{-1}(\Omega; \mathbb{R}^d)$ and $g\in L^2_\chi(\Omega)$. Let $(\boldsymbol\phi, p)\in H_0^2(\Omega; \mathbb R^d)\times(H_0^1(\Omega)\cap L_0^2(\Omega))$ be the solution of problem~\eqref{eq:fourthorderStokes}. Assume $\boldsymbol{\phi}\in H^3(\Omega; \mathbb R^d)$ and $p\in H^2(\Omega)$ which admit the estimate
\begin{equation}\label{eq:fourthorderStokesregularity1}
\|\boldsymbol\phi\|_3+\|p\|_2\lesssim \|\boldsymbol f\|_{-1} + \|g\|_{0}.
\end{equation}
\end{assumption}


\begin{assumption}\label{RegularityAssumption2}
	Assume that $\boldsymbol f\in H^{-1}(\Omega; \mathbb R^d)$. Let $\boldsymbol u\in H_0^2(\Omega; \mathbb R^d)$ be the solution of the following problem:
	\begin{equation*}
	\begin{cases}
	\Delta (\div\boldsymbol\sigma(\boldsymbol u))=\boldsymbol f & {\rm  in }\; \Omega,\\
	\boldsymbol u=\partial_{\boldsymbol{n}}  \boldsymbol u= \boldsymbol 0 & {\rm on }\; \partial\Omega,
	\end{cases}
	\end{equation*}
	where $\lambda \in[0,\Lambda]$ and $\mu\in [\mu_0,\mu_1]$ being the Lam\'{e} constants. Then $\boldsymbol u\in H^3(\Omega; \mathbb R^d)$ and it admits the estimate
	\[
	\| \boldsymbol u\|_3 \lesssim \|\boldsymbol f\|_{-1},
	\]
	where the hidden constant may depend on $\Lambda$, $\mu_0$ and $\mu_1$.
\end{assumption}

% We will prove Assumption~\ref{RegularityAssumption1} holds for a convex polygon in Section~\ref{sec:assumption12d}. 
Since the solution $\boldsymbol{\phi}$ of \eqref{eq:fourthorderStokes} is zero when $\boldsymbol{f}\in\nabla L^2(\Omega)$,
by replacing $\boldsymbol{f}$ in problem \eqref{eq:fourthorderStokes} with $\boldsymbol{f}+\nabla q$, we have from \eqref{eq:fourthorderStokesregularity1} that
\begin{equation}\label{eq:fourthorderStokesregularity2}
\|\boldsymbol\phi\|_3\lesssim
\inf_{q\in L^2(\Omega)}\|\boldsymbol f+\nabla q\|_{-1}.
\end{equation}
 
The Assumption \ref{RegularityAssumption1} with $g=0$ is exactly the Assumption 3.1 in \cite{ChenHuangHuang2023}, which has been proved for convex polygons in two dimensions in \cite[Theorem 3.6]{ChenHuangHuang2023} by using the $H^4$ regularity result for the triharmonic equation.
Then we prove Assumption \ref{RegularityAssumption1} for general $g$ in two dimensions for convex polygons in the remainder of this section.


% \subsection{Verification of Assumptions in two dimensions}

\begin{lemma}\label{lem:main}
Let $\Omega\subset\mathbb R^2$ be a bounded, simply connected polygonal domain.
Then
\[
\div\big( H^3(\Omega;\mathbb R^2)\cap  H_0^2(\Omega;\mathbb R^2)\big)
=
 H^2(\Omega)\cap  H_0^1(\Omega)\cap L_0^2(\Omega).
\]
Moreover, there exists a bounded linear right inverse
\[
\mathcal R:\  H^2(\Omega)\cap  H_0^1(\Omega)\cap L_0^2(\Omega)
\to
 H^3(\Omega;\mathbb R^2)\cap  H_0^2(\Omega;\mathbb R^2),
\quad
\div(\mathcal R g)=g.
\]
\end{lemma}

\begin{proof}
It can be verified directly that 
\[
\div\big( H^3(\Omega;\mathbb R^2)\cap  H_0^2(\Omega;\mathbb R^2)\big)
\subseteq
 H^2(\Omega)\cap  H_0^1(\Omega)\cap L_0^2(\Omega).
\]
Then we prove the reverse inclusion and the existence of the bounded linear right inverse.
Let
$
g\in  H^2(\Omega)\cap  H_0^1(\Omega)\cap L_0^2(\Omega).
$
We shall construct $\boldsymbol u\in  H^3(\Omega;\mathbb R^2)\cap  H_0^2(\Omega;\mathbb R^2)$ with $\div\boldsymbol u=g$.

Noting that $g|_{\partial\Omega}=0$, by Theorem 3.1 in \cite{ArnoldScottVogelius1988}, there exists a $\boldsymbol u_1\in H^3(\Omega;\mathbb R^2)\cap  H_0^1(\Omega;\mathbb R^2)$ such that
\[
\div\boldsymbol u_1 = g,
\qquad
\|\boldsymbol u_1\|_{3}\lesssim \|g\|_{2}.
\]
Let $\Delta_0(\Omega)$ and $\Delta_1(\Omega)$ be set of all vertices and edges of the polygon $\Omega$, respectively. For each edge $e\in \Delta_1(\Omega)$, let $\boldsymbol n$ and $\boldsymbol t$ be the unit outward normal and tangential vectors on $e$, respectively.
By $\boldsymbol u_1|_{\partial\Omega}=0$ and $g|_{\partial\Omega}=0$, we have
% Fix an open edge $e\subset\partial\Omega$.
% Since $\boldsymbol{n},\boldsymbol{\tau}$ are constant on $e$ and $\boldsymbol u_1|_e=0$, we have $\partial_\tau \boldsymbol u_1|_e=0$ (trace sense).
% Moreover, $g\in  H_0^1(\Omega)$ implies $g|_{\partial\Omega}=0$, hence $(\div\boldsymbol u_1)|_e=g|_e=0$.
% On a straight edge,
% \[
% (\div\boldsymbol u_1)|_e
% =
% (\boldsymbol{n}\cdot \partial_n\boldsymbol u_1)|_e + (\boldsymbol{t}\cdot \partial_t\boldsymbol u_1)|_e
% =
% (\boldsymbol{n}\cdot \partial_n\boldsymbol u_1)|_e,
% \]
% so
\begin{equation}\label{eq:normal-comp-zero}
(\boldsymbol{n}\cdot \partial_n\boldsymbol u_1)|_e = (\boldsymbol{n}\cdot \partial_n\boldsymbol u_1)|_e + (\boldsymbol{t}\cdot \partial_t\boldsymbol u_1)|_e
= (\div\boldsymbol u_1)|_e = 0.
\end{equation}
% Thus the boundary normal derivative $\partial_n\boldsymbol u_1$ is purely tangential on each edge.
% \smallskip
% \noindent\textbf{2.3 Define the tangential mismatch $h$.}
Define the scalar boundary function
\[
h := \boldsymbol{t}\cdot (\partial_n\boldsymbol u_1)\big|_{\partial\Omega}=\partial_n\boldsymbol u_1|_{\partial\Omega}.
\]
Since $\boldsymbol u_1\in  H^3(\Omega;\mathbb R^2)$, we have $\partial_n\boldsymbol u_1|_{\partial\Omega}\in  H^{3/2}(\partial\Omega;\mathbb R^2)$ and hence
\begin{equation}\label{eq:h-reg}
h\in  H^{3/2}(\partial\Omega),
\qquad
\|h\|_{ H^{3/2}(\partial\Omega)}\lesssim \|\boldsymbol u_1\|_{3}.
\end{equation}
For vertex $\texttt{v}\in\Delta_0(\Omega)$ shared by two edges $e^\pm \in \Delta_1(\Omega)$. Because $\boldsymbol u_1\in  H^3$ and $\boldsymbol u_1|_{\partial\Omega}=0$ as a global trace, we have $\partial_t\boldsymbol u_1|_{e^\pm}=0$. This means that $(\nabla \boldsymbol{u}_1)(\texttt{v})=0$. Hence
\begin{equation}\label{eq:vertex-normal-derivative}
% (\partial_n \boldsymbol u_1)(\texttt{v})=\boldsymbol 0,
% \qquad
% \text{hence}\qquad
h(\texttt{v})=0
\ \ \text{for each vertex }\texttt{v}\in\Delta_0(\Omega).
\end{equation}
Then there exists a $\psi\in H^4(\Omega)\cap H_0^2(\Omega)$ such that (cf. \cite[Theorem~5]{Agranovich2007}, \cite[Theorem~4]{Agranovich2008} and \cite[Theorem~R(m)]{Verchota1990})
\begin{equation}\label{eq:psi-traces}
\partial_{nn}\psi|_{\partial\Omega}=-h,\quad \textrm{ and } \quad \|\psi\|_{ H^4(\Omega)}\lesssim \|h\|_{ H^{3/2}(\partial\Omega)}.
\end{equation}

Set
$\boldsymbol v := \curl\psi \in  H^3(\Omega;\mathbb R^2)\cap H_0^1(\Omega;\mathbb R^2)$, then
$\div\boldsymbol v = 0$.
We now verify that $\boldsymbol v$ exactly cancels the boundary normal derivative of $\boldsymbol u_1$.
On each edge $e\in\Delta_1(\Omega)$,
\[
\boldsymbol{t}\cdot(\partial_n\boldsymbol v)
=
\partial_n(\boldsymbol v\cdot \boldsymbol{t})
=
\partial_n(-\partial_n\psi)
=
-\partial_{nn}\psi
=
h
=
\boldsymbol{t}\cdot(\partial_n\boldsymbol u_1),
\]
and
\[
\boldsymbol{n}\cdot(\partial_n\boldsymbol v)
=
\partial_n(\boldsymbol v\cdot \boldsymbol{n})
=
\partial_n(\partial_t\psi)
=
\partial_t(\partial_n\psi)
=
0.
\]
% where the last step uses $\partial_n\psi|_{\partial\Omega}=0$.
Together with \eqref{eq:normal-comp-zero}, this yields
\begin{equation}\label{eq:normal-derivative-match}
(\partial_n\boldsymbol v)|_{\partial\Omega}=(\partial_n\boldsymbol u_1)|_{\partial\Omega}.
\end{equation}
Let
$\boldsymbol u := \boldsymbol u_1 - \boldsymbol v$.
Then $\boldsymbol u\in  H^3(\Omega;\mathbb R^2)\cap H_0^1(\Omega;\mathbb R^2)$ and
\[
\div\boldsymbol u = \div\boldsymbol u_1 - \div\boldsymbol v = g - 0 = g.
\]
From \eqref{eq:normal-derivative-match} we obtain $\partial_n\boldsymbol u|_{\partial\Omega}=0$.
Thus $\boldsymbol u\in  H^3(\Omega;\mathbb R^2)\cap  H_0^2(\Omega;\mathbb R^2)$ and $\div\boldsymbol u=g$.

Moreover, the construction is linear in $g$ and the estimates
\[
\|\boldsymbol u_1\|_{3}\lesssim \|g\|_{ 2},
\quad
\|h\|_{ H^{3/2}(\partial\Omega)}\lesssim \|\boldsymbol u_1\|_{3},
\quad
\|\psi\|_{4}\lesssim \|h\|_{ H^{3/2}(\partial\Omega)},
\quad
\|\boldsymbol v\|_{3}\lesssim \|\psi\|_{4}
\]
imply
$
\|\boldsymbol u\|_{3}\lesssim \|g\|_{2}.
$
\end{proof}



\section{Regularity results}\label{sec:regularity}

Now we are ready to show the regularity results of the SGE model \eqref{eq:straingradproblem} under Assumptions~\ref{RegularityAssumption0}-\ref{RegularityAssumption2}.


% \subsection{Regularity of the SGE model}
% Now we are ready to show the regularity results of the SGE model. From now on we always suppose Assumptions~\ref{RegularityAssumption1}-\ref{RegularityAssumption2} hold.
\begin{lemma}\label{lm:fourthorderelasregularity}
	Let $\boldsymbol w\in H_0^2(\Omega; \mathbb R^d)$ be the solution of
	\begin{equation}\label{eq:fourthorderelasproblem}
	\begin{cases}
	\div\Delta\big(2\mu\boldsymbol\varepsilon(\boldsymbol w)+\lambda(\div\boldsymbol w)\boldsymbol I\big)=\boldsymbol f & {\rm  in }\; \Omega,\\
	\boldsymbol w=\partial_{\boldsymbol{n}}\boldsymbol w=\boldsymbol 0 & {\rm on }\; \partial\Omega
	\end{cases}
	\end{equation}
with $\boldsymbol f\in H^{-1}(\Omega; \mathbb{R}^d)$.
	Suppose Assumptions \ref{RegularityAssumption1}-\ref{RegularityAssumption2} hold.
	We have
	\begin{align}\label{eq:fourthorderelasregularity}
	\|\boldsymbol w\|_{2+j}+\lambda\|\div\boldsymbol w\|_{1+j}\lesssim \|\boldsymbol f\|_{j-2} \quad\textrm{ for } j=0,1.
	\end{align}
The hidden constants may depend on $\Omega$, $\mu$ and the constants
in Assumptions \ref{RegularityAssumption1}--\ref{RegularityAssumption2},
but are independent of $\lambda$.
\end{lemma}
\begin{proof}
The estimate \eqref{eq:fourthorderelasregularity} for $j=0$ has been proved in \cite[Lemma 3.7]{ChenHuangHuang2023}. 
Then we prove \eqref{eq:fourthorderelasregularity} for $j=1$. 
	% We follow the argument of Lemma~2.2 in \cite{BrennerSung1992} to prove \eqref{eq:fourthorderelasregularity}.
	% The weak formulation of \eqref{eq:fourthorderelasproblem} is
	% \begin{equation}\label{eq:20210918}
	% 2\mu(\nabla\boldsymbol\varepsilon(\boldsymbol w), \nabla\boldsymbol\varepsilon(\boldsymbol v)) + \lambda(\nabla\div\boldsymbol w, \nabla\div\boldsymbol v)=(\boldsymbol f, \boldsymbol v) \quad\forall~\boldsymbol v\in H_0^2(\Omega; \mathbb R^d),
	% \end{equation}%\cite{LaxMilgram1954}
	% which is well-posed by the $H^2$-Korn's inequality \eqref{eq:korninequalityH2}, and it holds
	% $
	% \|\boldsymbol w\|_2\lesssim \|\boldsymbol f\|_{-2}.
	% $
	% By Lemma~\ref{lm:divrightinverse}, there exists $\widetilde{\boldsymbol w}\in H_0^2(\Omega; \mathbb R^d)$ such that
	% $\div\widetilde{\boldsymbol w}=\div\boldsymbol w$ and $\|\widetilde{\boldsymbol w}\|_2\lesssim \|\div\boldsymbol w\|_1.$
	% Taking $\boldsymbol v=\widetilde{\boldsymbol w}$ in \eqref{eq:20210918}, it follows that
	% \begin{equation*}
	% \lambda|\div\boldsymbol w|_1^2=(\boldsymbol f, \widetilde{\boldsymbol w})-2\mu(\nabla\boldsymbol\varepsilon(\boldsymbol w), \nabla\boldsymbol\varepsilon(\widetilde{\boldsymbol w}))\lesssim \|\boldsymbol f\|_{-2}\|\widetilde{\boldsymbol w}\|_2+\|\boldsymbol w\|_2\|\widetilde{\boldsymbol w}\|_2,
	% \end{equation*}
	% which gives \eqref{eq:fourthorderelasregularity} for $j=0$.

Let $(\boldsymbol\phi,p_0)\in H_0^2(\Omega;\mathbb R^d)\times
\big(H_0^1(\Omega)\cap L_0^2(\Omega)\big)$ solve
\begin{equation}\label{eq:limit_stokes_journal2}
\begin{cases}
\Delta^2\boldsymbol\phi+\nabla\Delta p_0=\dfrac{1}{\mu}\boldsymbol f
& \text{in }\Omega,\\[2pt]
\Delta\operatorname{div}\boldsymbol\phi=0 & \text{in }\Omega.
\end{cases}
\end{equation}
By Assumption~\ref{RegularityAssumption1},
\begin{equation}\label{eq:stokes_reg_journal2}
\|\boldsymbol\phi\|_{3}+\|p_0\|_{2}\lesssim\|\boldsymbol f\|_{-1}.
\end{equation}
Introduce
\[
% \pi:=(\mu+\lambda)\operatorname{div}\boldsymbol w,\qquad
\gamma:=\frac{\mu}{\mu+\lambda}\in(0,1],\qquad
\eta:=\frac{1}{\gamma}\operatorname{div}\boldsymbol w-p_0,\qquad \boldsymbol e:=\boldsymbol w-\boldsymbol\phi.
\]
The first equation in \eqref{eq:fourthorderelasproblem} can be rewritten as
\begin{equation}\label{eq:elas_rewrite_journal2}
\Delta^2\boldsymbol w+\nabla\Delta\Big(\frac{1}{\gamma}\operatorname{div}\boldsymbol w\Big)=\frac{1}{\mu}\boldsymbol f  \quad \textrm{ in } \Omega.
\end{equation}
Subtracting \eqref{eq:limit_stokes_journal2} gives
\begin{equation}\label{eq:error_system_journal2}
\begin{cases}
\Delta^2\boldsymbol e+\nabla\Delta\eta=\boldsymbol 0 & \text{in }\Omega,\\[2pt]
\Delta\operatorname{div}\boldsymbol e-\gamma\,\Delta\eta
=\gamma\,\Delta p_0& \text{in }\Omega,\\[2pt]
\boldsymbol e=\partial_{\boldsymbol n}\boldsymbol e=\boldsymbol 0 &\text{on }\partial\Omega.
\end{cases}
\end{equation}

We now use a standard perturbation argument. By Lemma~\ref{lem:L2chi}, $\Delta p_0\in L^2_\chi(\Omega)$ and hence $\varepsilon\,\Delta p_0\in L^2_\chi(\Omega)$.
Set
\begin{align*}
X&:=\big(H^3(\Omega;\mathbb R^d)\cap H_0^2(\Omega;\mathbb R^d)\big)\times
\big(H^2(\Omega)\cap H_0^1(\Omega)\cap L_0^2(\Omega)\big), \\
Y&:=H^{-1}(\Omega;\mathbb R^d)\times L^2_\chi(\Omega),
\end{align*}
and
define operators $\mathcal A_0,\mathcal B\in\mathcal L(X,Y)$ by
\[
\mathcal A_0(\boldsymbol v,r):=\big(\Delta^2\boldsymbol v+\nabla\Delta r,\ \Delta\operatorname{div}\boldsymbol v\big),
\qquad
\mathcal B(\boldsymbol v,r):=\big(\boldsymbol 0,\ \Delta r\big).
\]
The norms for spaces $X$ and $Y$ are given by $$
\|(\boldsymbol v,r)\|_X^2:=\|\boldsymbol v\|_3^2+\|r\|_2^2,
\qquad
\|(\boldsymbol v,r)\|_Y^2:=\|\boldsymbol v\|_{-1}^2+\|r\|_0^2.
$$	
Assumption~\ref{RegularityAssumption1} implies that $\mathcal A_0:X\to Y$ is an isomorphism and
\begin{equation}\label{eq:A0inv_bound_journal2}
\|(\boldsymbol v,r)\|_{X}\le C_S\,\|\mathcal A_0(\boldsymbol v,r)\|_{Y}
\qquad\forall\,(\boldsymbol v,r)\in X.
\end{equation}
We can choose $C_S\geq1$.
Moreover, $\|\mathcal B(\boldsymbol v,r)\|_Y=\|\Delta r\|_0\le \sqrt{d}\|r\|_2\le \sqrt{d}\|(\boldsymbol v,r)\|_X$,
hence $\|\mathcal B\|_{\mathcal L(X,Y)}\le \sqrt{d}$.
Consequently,
\[
\|\mathcal A_0^{-1}\mathcal B\|_{\mathcal L(X)}\le
\|\mathcal A_0^{-1}\|_{\mathcal L(Y,X)}\,\|\mathcal B\|_{\mathcal L(X,Y)}
\le \sqrt{d}\,C_S.
\]
Choose
$
\gamma_0:=\frac{1}{2\sqrt{d}\,C_S}\in(0,1).
$
If $\gamma\le\gamma_0$, equivalently 
$$
\lambda\ge \lambda_\ast:=\mu\Big(\gamma_0^{-1}-1\Big)=\mu(2\sqrt{d}\,C_S-1),$$
then the operator
$\mathcal A_\gamma:=\mathcal A_0-\gamma\mathcal B$ is invertible on $X$ and
\begin{equation}\label{eq:perturb_inverse_journal2}
\|(\boldsymbol v,r)\|_X\le 2C_S\,\|\mathcal A_\gamma(\boldsymbol v,r)\|_Y
\qquad\forall\,(\boldsymbol v,r)\in X,
\end{equation}
which follows from the Neumann series for
$I-\gamma\mathcal A_0^{-1}\mathcal B$ (cf. \cite[Chapter~I, \S4, Example~4.5, (4.22) and the perturbation estimate (4.24), pp.~30--31]{Kato1995}).
Applying \eqref{eq:perturb_inverse_journal2} to \eqref{eq:error_system_journal2} gives
\begin{equation}\label{eq:error_est_journal2}
\|\boldsymbol e\|_3+\|\eta\|_2\le 2C_S\,\gamma\,\|\Delta p_0\|_0
\le 2\sqrt{d}\,C_S\,\gamma\,\|p_0\|_2\leq \|p_0\|_2\quad\textrm{ for }\lambda\ge\lambda_\ast.
\end{equation}
Combining \eqref{eq:stokes_reg_journal2} and \eqref{eq:error_est_journal2} yields
\[
\|\boldsymbol w\|_3\le \|\boldsymbol\phi\|_3+\|\boldsymbol e\|_3
\lesssim\|\boldsymbol f\|_{-1}\quad\textrm{ for }\lambda\ge\lambda_\ast.
\]
For the divergence term, since $\operatorname{div}\boldsymbol w=\gamma(\eta+p_0)$,
\[
(\mu+\lambda)\|\operatorname{div}\boldsymbol w\|_2
\le \mu(\|\eta\|_2+\|p_0\|_2)
\le 2\mu\,\|p_0\|_2\lesssim \|\boldsymbol f\|_{-1},
\]
and therefore $\lambda\|\operatorname{div}\boldsymbol w\|_2\lesssim \|\boldsymbol f\|_{-1}$ as well.
Hence \eqref{eq:fourthorderelasregularity} holds for all $\lambda\ge\lambda_\ast$ with a constant independent of $\lambda$.

For $0\le\lambda\le\lambda_\ast$, \eqref{eq:elas_rewrite_journal2} is precisely
$\Delta(\div\boldsymbol\sigma(\boldsymbol w))=\boldsymbol f$.
Assumption~\ref{RegularityAssumption2} yields
\[
\|\boldsymbol w\|_3\lesssim \|\boldsymbol f\|_{-1},
\qquad
\lambda\|\operatorname{div}\boldsymbol w\|_2\le \lambda_\ast\|\operatorname{div}\boldsymbol w\|_2
\lesssim \lambda_\ast\|\boldsymbol w\|_3\lesssim \|\boldsymbol f\|_{-1}.
\]
Together with the large--$\lambda$ estimate, this proves \eqref{eq:fourthorderelasregularity} for all $\lambda\ge0$.
\end{proof}

\begin{lemma}
Suppose Assumptions \ref{RegularityAssumption0}-\ref{RegularityAssumption2} hold.
Let $\boldsymbol w\in H_0^2(\Omega; \mathbb R^d)$ be the solution of equation \eqref{eq:fourthorderelasproblem}.
We have
\begin{align}
\label{eq:fourthorderelasregularityw}
\|\boldsymbol w\|_{3}\lesssim \inf_{q\in L^2(\Omega)}\|\boldsymbol f+\nabla q\|_{-1} + \frac{1}{\lambda}\|\boldsymbol f\|_{-1}.
\end{align}
\end{lemma}
\begin{proof}
Take any $q\in L^2(\Omega)$. Let $C_q=(q,\chi)/(1,\chi)\in\mathbb R$, then $q-C_q\in L^2_\chi(\Omega)$. By \eqref{eq:Deltasurjection}, there exists $p\in H^2(\Omega)\cap H_0^1(\Omega)\cap L_0^2(\Omega)$ such that $\Delta p=q-C_q$, which indicates $\nabla\Delta p=\nabla q$.
Thus it holds from \eqref{eq:limit_stokes_journal2} that
\[
\Delta^2\boldsymbol\phi+\nabla\Delta\Big(p_0+\frac{1}{\mu}p\Big)
=\frac{1}{\mu}\boldsymbol f+\frac{1}{\mu}\nabla q.
\]
Apply the Assumption~\ref{RegularityAssumption1} to acquire
\[
\|\boldsymbol\phi\|_3\lesssim \|\boldsymbol f+\nabla q\|_{-1}\qquad\forall\,q\in L^2(\Omega),
\]
hence
\begin{equation}\label{eq:phi_bestgrad_journal2}
\|\boldsymbol\phi\|_3\lesssim \inf_{q\in L^2(\Omega)}\|\boldsymbol f+\nabla q\|_{-1}.
\end{equation}

If $\lambda\ge\lambda_\ast$, then \eqref{eq:error_est_journal2} and \eqref{eq:stokes_reg_journal2} imply
\[
\|\boldsymbol e\|_3\lesssim \gamma\,\|p_0\|_2\lesssim  \frac{\mu_1}{\lambda}\,\|\boldsymbol f\|_{-1}\qquad(\lambda>0).
\]
Combining with \eqref{eq:phi_bestgrad_journal2} yields \eqref{eq:fourthorderelasregularityw} for $\lambda\ge\lambda_\ast$.

For $0<\lambda<\lambda_\ast$, Assumption~\ref{RegularityAssumption2} gives
$\|\boldsymbol w\|_3\lesssim \|\boldsymbol f\|_{-1}\le (\lambda_\ast/\lambda)\|\boldsymbol f\|_{-1}$.
Since $\inf_{q\in L^2}\|\boldsymbol f+\nabla q\|_{-1}\ge0$, we obtain
\[
\|\boldsymbol w\|_3
\lesssim \inf_{q\in L^2(\Omega)}\|\boldsymbol f+\nabla q\|_{-1}+\frac{1}{\lambda}\|\boldsymbol f\|_{-1},
\]
after absorbing $\lambda_\ast$ into the hidden constant. This completes the proof.	
\end{proof}

Taking $\iota=0$, problem \eqref{eq:straingradproblem} becomes the linear elasticity model. Let $\boldsymbol u_0\in H_0^1(\Omega; \mathbb R^d)$ be the solution of the linear elasticity problem
\begin{equation}\label{eq:elasticity}
\begin{cases}
-\mu\Delta\boldsymbol u_0 - (\lambda+\mu)\nabla\div\boldsymbol u_0 =\boldsymbol f & \textrm{ in } \Omega,\\
\boldsymbol u_0=\boldsymbol0 & \textrm{ on } \partial\Omega.
\end{cases}
\end{equation}
According to \cite{BrennerSung1992,MazyaRossmann2010}, it holds
\begin{equation}\label{eq:u0regularity}
\|\boldsymbol u_0\|_2+\lambda\|\div\boldsymbol u_0\|_1 \lesssim \|\boldsymbol f\|_0.
\end{equation}


\begin{theorem}\label{tm:upregularity}
Suppose Assumptions \ref{RegularityAssumption0}-\ref{RegularityAssumption2} hold.
Let $\boldsymbol u\in H_0^2(\Omega; \mathbb R^d)$ be the solution of problem \eqref{eq:straingradproblem}, and $\boldsymbol u_0\in H_0^1(\Omega; \mathbb R^d)$ satisfy the linear elasticity problem \eqref{eq:elasticity}. It holds
\begin{equation}\label{eq:straingradregularity}	
\begin{aligned}	
&|\boldsymbol u-\boldsymbol u_0|_{1} +\iota\|\boldsymbol u\|_2  + \iota^{2}\|\boldsymbol u\|_{3}\\
&\qquad+\lambda\|\div(\boldsymbol u-\boldsymbol u_0)\|_0 + \lambda\iota|\div\boldsymbol u|_1 + \lambda\iota^2\|\div\boldsymbol u\|_2 \lesssim \iota^{1/2}\|\boldsymbol f\|_0.
\end{aligned}
\end{equation}
\end{theorem}
\begin{proof}
Recall the following estimate from \cite[Lemma 3.8]{ChenHuangHuang2023}:
\begin{equation}\label{eq:straingradregularity1}
	|\boldsymbol u-\boldsymbol u_0|_{1} +\iota\|\boldsymbol u\|_2  + \iota^{2}\|\boldsymbol u\|_{3} \lesssim \iota^{1/2}\|\boldsymbol f\|_0.
	\end{equation}
Applying Lemma~\ref{lm:fourthorderelasregularity} to %\eqref{eq:20210919}, 
	\begin{equation}\label{eq:20210919}
	\div\Delta\boldsymbol\sigma(\boldsymbol u)=\iota^{-2}\div\boldsymbol\sigma(\boldsymbol u-\boldsymbol u_0),
	\end{equation}
we obtain from \eqref{eq:fourthorderelasregularity} and \eqref{eq:straingradregularity1} that
	\begin{align}
	\lambda\iota^{2}\|\div\boldsymbol u\|_{2}&\lesssim \|\div\boldsymbol\sigma(\boldsymbol u-\boldsymbol u_0)\|_{-1}\lesssim \|\boldsymbol\sigma(\boldsymbol u-\boldsymbol u_0)\|_{0}\lesssim |\boldsymbol u-\boldsymbol u_0|_1
	+\lambda\|\div(\boldsymbol u-\boldsymbol u_0)\|_0 \notag\\
	&\lesssim \iota^{1/2}\|\boldsymbol f\|_0+\lambda\|\div(\boldsymbol u-\boldsymbol u_0)\|_0. \label{eq:20210923-4}
	\end{align}
Thanks to \eqref{eq:straingradregularity1} and \eqref{eq:20210923-4}, it suffices to prove
	\begin{equation}\label{eq:straingradregularity0}
	\lambda\|\div(\boldsymbol u-\boldsymbol u_0)\|_0 + \lambda\iota|\div\boldsymbol u|_1 \lesssim \iota^{1/2}\|\boldsymbol f\|_0.
	\end{equation}

By Theorem~1.1 in \cite{CostabelMcIntosh2010}, there exists $\boldsymbol v\in H_0^1(\Omega; \mathbb R^d)$ such that
	\begin{equation}\label{eq:20210923-3}
	\div\boldsymbol v=\div(\boldsymbol u-\boldsymbol u_0),\quad \|\boldsymbol v\|_1\lesssim \|\div(\boldsymbol u-\boldsymbol u_0)\|_0.
	\end{equation}
	Multiply \eqref{eq:20210919} by $\boldsymbol v$ and apply the integration by parts to get
	\begin{align}
	&\quad \lambda\|\div(\boldsymbol u-\boldsymbol u_0)\|_0^2 +\lambda\iota^{2}|\div\boldsymbol u|_1^2 \notag\\
	&=-2\mu(\boldsymbol\varepsilon(\boldsymbol u-\boldsymbol u_0), \boldsymbol\varepsilon(\boldsymbol v)) + 2\mu\iota^{2}(\Delta\boldsymbol\varepsilon(\boldsymbol u), \boldsymbol\varepsilon(\boldsymbol v)) - \lambda\iota^{2}(\Delta\div\boldsymbol u, \div\boldsymbol u_0). \label{eq:20211217}
	\end{align}
	By the multiplicative trace inequality and \eqref{eq:20210923-4}, it holds
	\begin{align*}
	- \lambda\iota^{2}(\Delta\div\boldsymbol u, \div\boldsymbol u_0)&=\lambda\iota^{2}(\nabla\div\boldsymbol u, \nabla\div\boldsymbol u_0)- \lambda\iota^{2}(\partial_{\boldsymbol{n}}\div\boldsymbol u, \div\boldsymbol u_0)_{\partial\Omega} \\
	&\lesssim \lambda\iota^{2}|\div\boldsymbol u|_1|\div\boldsymbol u_0|_1 + \lambda\iota^{2}|\div\boldsymbol u|_1^{1/2}\|\div\boldsymbol u\|_2^{1/2}\|\div\boldsymbol u_0\|_1 \\
	&\lesssim \lambda\iota^{2}|\div\boldsymbol u|_1|\div\boldsymbol u_0|_1 + \lambda^{1/2}\iota^{5/4}|\div\boldsymbol u|_1^{1/2}\|\boldsymbol f\|_0^{1/2}\|\div\boldsymbol u_0\|_1 \\
	&\quad + \lambda\iota|\div\boldsymbol u|_1^{1/2}\|\div(\boldsymbol u-\boldsymbol u_0)\|_0^{1/2}\|\div\boldsymbol u_0\|_1.
	\end{align*}
	Then we acquire from \eqref{eq:20210923-3}-\eqref{eq:20211217} that
	\begin{align*}
	&\quad \lambda\|\div(\boldsymbol u-\boldsymbol u_0)\|_0^2 +\lambda\iota^{2}|\div\boldsymbol u|_1^2 \\
	&\lesssim (|\boldsymbol u-\boldsymbol u_0|_1 + \iota^{2}|\boldsymbol u|_3)|\boldsymbol v|_1 +\lambda\iota^{2}|\div\boldsymbol u|_1|\div\boldsymbol u_0|_1 \\
	&\quad + \lambda^{1/2}\iota^{5/4}|\div\boldsymbol u|_1^{1/2}\|\boldsymbol f\|_0^{1/2}\|\div\boldsymbol u_0\|_1+ \lambda\iota|\div\boldsymbol u|_1^{1/2}\|\div(\boldsymbol u-\boldsymbol u_0)\|_0^{1/2}\|\div\boldsymbol u_0\|_1 \\
	&\lesssim (|\boldsymbol u-\boldsymbol u_0|_1 + \iota^{2}|\boldsymbol u|_3) \|\div(\boldsymbol u-\boldsymbol u_0)\|_0 +(\lambda\iota^{2}|\div\boldsymbol u|_1^2)^{1/2}(\lambda\iota^2\|\div\boldsymbol u_0\|_1^2)^{1/2} \\
	&\quad + \iota^{3/2}|\div\boldsymbol u|_1\|\boldsymbol f\|_0 + \lambda\iota\|\div\boldsymbol u_0\|_1^2 \\
	&\quad + (\lambda\iota|\div\boldsymbol u|_1\|\div(\boldsymbol u-\boldsymbol u_0)\|_0)^{1/2}(\lambda\iota\|\div\boldsymbol u_0\|_1^2)^{1/2}.
	\end{align*}
	This implies
	\begin{align*}
	&\quad \lambda\|\div(\boldsymbol u-\boldsymbol u_0)\|_0^2 +\lambda\iota^{2}|\div\boldsymbol u|_1^2 \lesssim \frac{1}{\lambda}(|\boldsymbol u-\boldsymbol u_0|_1 + \iota^{2}|\boldsymbol u|_3)^2+ \lambda\iota\|\div\boldsymbol u_0\|_1^2 + \frac{1}{\lambda}\iota\|\boldsymbol f\|_0^2,
	\end{align*}
	i.e.,
	\begin{equation*}
	\lambda\|\div(\boldsymbol u-\boldsymbol u_0)\|_0 +\lambda\iota|\div\boldsymbol u|_1 \lesssim |\boldsymbol u-\boldsymbol u_0|_1 + \iota^{2}|\boldsymbol u|_3+\lambda\iota^{1/2}\|\div\boldsymbol u_0\|_1+\iota^{1/2}\|\boldsymbol f\|_0.
	\end{equation*}
	Therefore, we conclude \eqref{eq:straingradregularity0} from
	\eqref{eq:straingradregularity1} and \eqref{eq:u0regularity}.
\end{proof}



\bibliographystyle{abbrv}
\bibliography{refs}

\end{document}
